% 本文件由 LaTeX 排版 Agent 根据有效 translation.json 自动生成。
% author=Ignatius of Antioch; work_id=0106; title=依纳爵致特拉利人书

\chapter{依纳爵致特拉利人书}

% structure: p0001 source=h1 target=omitted reason=page-title-already-rendered-by-parent

% structure: p0002 source=h2 target=section reason=relative-heading-rank; base=h2

\section{问候}

依纳爵，亦称塞奥弗鲁斯，致在亚细亚的特拉利的圣教会，蒙天主，即耶稣基督的圣父，所爱，被选，堪当承受天主的恩宠，因耶稣基督的肉身、宝血和苦难而得享平安。祂是我们的希望，因我们藉着祂而复活。我以完备的教会精神及品格向你们致候，并愿你们满渥喜乐。\par

% structure: p0004 source=h2 target=section reason=relative-heading-rank; base=h2

\section{第一章 称赞他们的美德}

\Emph{我知道}你们在忍耐中怀有无可指摘且真诚的心志，这不仅表现在当前的行为，更是出于本性，正如你们的主教波利比乌斯已向我表明。他因天主及耶稣基督的旨意来到士每拿，并深切体会到我——这为基督耶稣所束缚的人——所拥有的喜乐，以致我在他身上看到了你们全体会众。因此，我藉着他领受了你们按天主所怀善意的见证，我喜乐地发现，你们正如我所知，是天主的追随者。\par

% structure: p0006 source=h2 target=section reason=relative-heading-rank; base=h2

\section{第二章 服从主教等}

因为，既然你们服从主教如同服从耶稣基督，你们在我眼中看来，生活不是按人的方式，而是按耶稣基督。祂为我们而死，为使你们因信祂的死而脱离死亡。因此，正如你们现在所做的，没有主教的许可，你们不可做任何事，并且也应服从司铎团，如同服从耶稣基督的宗徒，祂是我们的希望。若我们活在祂内，终将得以显明。同样，执事作为耶稣基督奥秘的仆役，应当在一切事上令众人喜悦。因为他们不是酒食的仆役，而是天主的教会的仆人。因此，他们务必避开一切可受指控的把柄，如同避火一般。\par

% structure: p0008 source=h2 target=section reason=relative-heading-rank; base=h2

\section{第三章 尊敬执事等}

同样，众人应尊敬执事，视之为耶稣基督的任命；尊敬主教如耶稣基督，即圣父之子；尊敬司铎如天主的议会，宗徒的集会。离了这些，就没有教会。关于这一切，我相信你们也持有同样的看法。因为我已经领受了你们爱德的明证，并且这明证仍与我同在，即体现在你们的主教身上。他的仪容极富教导意义，他的温良本身就是一种力量。我甚至认为，不虔敬的人也必尊敬他，因为他们也乐于看见我不顾惜自己。然而，当我获许可就此点写信时，我岂能如此自负，竟以判罪之身，如同宗徒一般向你们发号施令？\par

% structure: p0010 source=h2 target=section reason=relative-heading-rank; base=h2

\section{第四章 我需要谦卑}

我在天主方面拥有伟大的知识，但我克制自己，唯恐因自夸而丧亡。因为如今我更需要谨慎，不可留意那些吹捧我的人。因为那些以赞扬之辞对我说话的人，是在鞭策我。因为我确实渴望受苦，但不知自己是否配得。这渴望虽不为多人所知，却越发猛烈地攻袭我。因此，我需要温良，藉此这世界的首领被消灭。\par

% structure: p0012 source=h2 target=section reason=relative-heading-rank; base=h2

\section{第五章 我不教导深奥道理}

我岂不能写信与你们谈论天上的事？但我畏惧，恐怕伤害你们这些在基督里尚为婴孩的人。请在此事上原谅我，免得你们因不能领受而窒息。因为即使我这为基督被捆绑的人，也未能因此理解天上的事，天使的居所，他们在各自首领下的聚集，以及可见与不可见的事物。撇开这些深奥的题目，我在其他方面仍是一个学习者；因为我们仍缺少许多，以不致亏缺于天主。\par

% structure: p0014 source=h2 target=section reason=relative-heading-rank; base=h2

\section{第六章 远离异端者的毒害}

因此，我——其实不是我，乃是耶稣基督的爱——恳求你们，只使用基督徒的食粮，避免另一种草料；我指的是异端。因为那些热衷于此的人，将耶稣基督与自己的毒药混杂，说不值得相信的话，如同将致命的毒药掺入甜酒，不知情的人贪婪地饮下，以致命的快乐导致自己的死亡。\par

% structure: p0016 source=h2 target=section reason=relative-heading-rank; base=h2

\section{第七章 续论同一主题}

因此，你们要提防这样的人。如果你们不自高自大，并与我们的天主耶稣基督、主教以及宗徒的诫命保持紧密的共融，那么你们就能做到这一点。凡在祭台之内的是洁净的，凡在外的则不洁净。也就是说，凡脱离主教、司铎团和执事而行事的人，他良心不洁。\par

% structure: p0018 source=h2 target=section reason=relative-heading-rank; base=h2

\section{第八章 提防魔鬼的罗网}

并非我知道你们中间有这种事，而是因为我深爱你们，预见魔鬼的罗网，所以提醒你们。因此，穿上温良，在信德——即主的肉身——及爱德——即耶稣基督的血——中更新自己。你们中不可有人对邻人怀怨。不可给予外邦人机会，免得因少数愚昧之人，使全体信天主的人被毁谤。因为，\Quote{因他的虚妄而让我的名字在任何地方受亵渎的人有祸了。} \BibleRef{依 52:5}\par

% structure: p0020 source=h2 target=section reason=relative-heading-rank; base=h2

\section{第九章 引用基督的史实}

因此，若有人向你们谈论与耶稣基督相悖的事，你们要掩耳不听。祂是大卫的后裔，也是玛利亚所生。祂真实地诞生、饮食。祂在本丢彼拉多手下真实地受迫害。祂真实地被钉十字架，且真实地死去，在天上、地上和地下的众生物眼前。祂也真实地从死者中复活，圣父使祂活过来。同样，圣父也将藉基督耶稣使信祂的我们复活。离了祂，我们就没有真生命。\par

% structure: p0022 source=h2 target=section reason=relative-heading-rank; base=h2

\section{第十章 基督苦难的真实性}

但若像某些无神的人——即不信者——所说，祂只是貌似受苦（他们自己也仅是貌似存在），那么我为何被捆绑？为何渴望被抛给野兽？我岂不是枉然而死？我对主的十字架岂不犯了虚谎之罪？\par

% structure: p0024 source=h2 target=section reason=relative-heading-rank; base=h2

\section{第十一章 躲避幻象论者的致命错误}

因此，你们要逃避那些〔\OriginalTerm{撒殚}{Satan}的〕邪恶枝芽，它们结出致死之果；任何人若品尝了它，便立即死亡。因为这些人不是圣父的栽种。倘若他们是，就会显出是十字架的枝条，而他们的果实必是不朽坏的。藉着这十字架，祂通过祂的苦难召唤你们，如同召唤祂的肢体。所以，头不能没有肢体而独自生成；天主，即是〔\OriginalTerm{救主}{Saviour}〕本身，已应许了他们的合一。\par

% structure: p0026 source=h2 target=section reason=relative-heading-rank; base=h2

\section{第十二章 继续在合一与爱德中}

我从\Place{斯弥纳}向你们请安，同我在一起的还有天主的众教会，他们在肉身和灵魂上处处使我获得安慰。我为耶稣基督的缘故佩戴的锁链（祈祷使我得以到达天主）劝勉你们。你们要继续彼此和谐共处，并彼此代祷；因为合乎你们每一位，尤其是长老们，去慰藉主教，为光荣圣父、耶稣基督和宗徒们。我在爱德中恳求你们听我，免得我因写了这些而成为反对你们的证据。也请为我祈祷，我需要你们的爱德和天主的仁慈，好使我堪当承受我所预定的份，并且不被视为可弃绝的。\par

% structure: p0028 source=h2 target=section reason=relative-heading-rank; base=h2

\section{第十三章 结论}

\Person{斯弥纳人}和\Person{厄弗所人}的爱德向你们致敬。请在你们的祈祷中记念在\Place{叙利亚}的教会，我也是从那里不配得到我的称呼，我是他们中最末的一个。愿你们在耶稣基督内安好，继续服从主教，如同服从〔\OriginalTerm{天主的}{of God}〕命令，同样也服从长老团。并且，每个人要以不可分割的心彼此相爱。愿我的\OriginalTerm{灵魂}{spirit}因你们而圣化，不仅现在，也在我到达天主时。因为我目前仍处于危险中。但圣父在耶稣基督内是信实的，必会成就我及你们的祈求：愿你们在他内被发现无可指责。\par

\SourceRecord{Translated by Alexander Roberts and James Donaldson.；Ante-Nicene Fathers；Vol. 1.；Edited by Alexander Roberts, James Donaldson, and A. Cleveland Coxe.；Buffalo, NY: Christian Literature Publishing Co.,；1885.；<http://www.newadvent.org/fathers/0106.htm>.}
